
\begin{lemma}\label{lem:exists_mem_conv_tendsto}
Let $(X_n)_{n \in \mathbb{N}}$ be a sequence of measurable functions from $\mathcal{X}$ to $\mathbb{R}_{+,\infty}$ and let $Q \in \mathcal{P}(\mathcal{X})$.
There exists a sequence of measurable functions $\tilde{X}_n \in \mathrm{conv}(X_n, X_{n+1}, \ldots)$ and $X : \mathcal{X} \to \mathbb{R}_{+, \infty}$ such that $\tilde{X}_n \to X$ $Q$-almost surely.
\end{lemma}

\begin{proof}
Proved in the Brownian motion project.
\end{proof}


\begin{lemma}\label{lem:exists_integral_utility_eq_sup}
  \uses{def:eVar}
Let $\mathcal{Q} \subseteq \mathcal{P}(\mathcal{X})$ be a set of probability measures and $U : \mathbb{R}_{+,\infty} \to \overline{\mathbb{R}}$ a concave function which is bounded from above.
Then there exists $f^* \in \mathcal{E}_{\mathcal{Q}}$ such that
\begin{align*}
  \sup_{f \in \mathcal{E}_{\mathcal{Q}}} P[U(f)]
  &= P[U(f^*)]
  \: .
\end{align*}
\end{lemma}

The properties of $\mathcal{E}_{\mathcal{Q}}$ we use for the proof of this lemma are convexity and the fact that if $f_n \in \mathcal{E}_{\mathcal{Q}}$ then $\liminf_{n \to \infty} f_n \in \mathcal{E}_{\mathcal{Q}}$.

\begin{proof}
  \uses{lem:exists_mem_conv_tendsto}
Let $(f_n)_{n \in \mathbb{N}}$ be a sequence of e-variables in $\mathcal{E}_{\mathcal{Q}}$ such that $n \mapsto P[U(f_n)]$ is non-decreasing and $P[U(f_n)] \to \sup_{f \in \mathcal{E}_{\mathcal{Q}}} P[U(f)]$.
Let $\tilde{f}_n$ be a sequence of measurable functions and let $\tilde{f} : \mathcal{X} \to \mathbb{R}_{+, \infty}$, with $\tilde{f}_n \in \mathrm{conv}(f_n, f_{n+1}, \ldots)$ such that $\tilde{f}_n \to \tilde{f}$ $P$-almost surely.
That sequence exists by Lemma~\ref{lem:exists_mem_conv_tendsto}.
By convexity of $\mathcal{E}_{\mathcal{Q}}$, $\tilde{f}_n \in \mathcal{E}_{\mathcal{Q}}$ for all $n$.

Let $f^* = \liminf_{n \to \infty} \tilde{f}_n$. Note that $f^*$ is a measurable function such that $f^* = \tilde{f}$ $P$-almost surely.
We have $f^* \in \mathcal{E}_{\mathcal{Q}}$ by an application of Fatou's lemma: for $Q \in \mathcal{Q}$,
\begin{align*}
  Q[f^*]
  = Q[\liminf_{n \to \infty} \tilde{f}_n]
  \le \liminf_{n \to \infty} Q[\tilde{f}_n]
  \le 1
  \: .
\end{align*}

By concavity of $U$ and the fact that $P[U(f_n)]$ is non-decreasing, we have
\begin{align*}
  P[U(\tilde{f}_n)]
  &= P[U(\sum_{i=n}^{m_n} \lambda_{n,i} f_{i})]
  \ge \sum_{i=n}^{m_n} \lambda_{n,i} P[U(f_{i})]
  \ge P[U(f_n)]
  \: .
\end{align*}

We also have $f^* = \limsup_{n \to \infty} \tilde{f}_n$ $P$-almost surely, hence by Fatou's lemma (using that $U$ is bounded from above),
\begin{align*}
  P[U(f^*)]
  &= P[U(\limsup_{n \to \infty} \tilde{f}_n)]
  \ge \limsup_{n \to \infty} P[U(\tilde{f}_n)]
  \ge \limsup_{n \to \infty} P[U(f_n)]
  = \sup_{f \in \mathcal{E}_{\mathcal{Q}}} P[U(f)]
  \: .
\end{align*}
Since $f^* \in \mathcal{E}_{\mathcal{Q}}$, we also have $P[U(f^*)] \le \sup_{f \in \mathcal{E}_{\mathcal{Q}}} P[U(f)]$, which proves the equality.
\end{proof}


\begin{lemma}\label{lem:exists_integral_utility_eq_sup_forall_lt_top}
  \uses{def:eVar}
Let $\mathcal{Q} \subseteq \mathcal{P}(\mathcal{X})$ be a set of probability measures and $U : \mathbb{R}_{+,\infty} \to \overline{\mathbb{R}}$ a concave function which is bounded from above.
Then there exists $f^* \in \mathcal{E}_{\mathcal{Q}}$ such that
\begin{align*}
  \sup_{f \in \mathcal{E}_{\mathcal{Q}}} P[U(f)]
  &= P[U(f^*)]
  \: ,
\end{align*}
and such that all e-variables are $P$-a.s. finite on $\{f^* < \infty\}$.
\end{lemma}

\begin{proof}
  \uses{lem:exists_integral_utility_eq_sup, lem:singularSet_unique, lem:nullSet_eq_top}
Let $f_0$ be as in Lemma~\ref{lem:exists_integral_utility_eq_sup}.
The set $\{f_0 = \infty\}$ is $\mathcal{Q}$-null by Lemma~\ref{lem:nullSet_eq_top}.

TODO
\end{proof}


\begin{lemma}\label{lem:bounded_numeraire_first_order}
  \uses{def:eVar, lem:exists_integral_utility_eq_sup_forall_lt_top}
Let $\mathcal{Q} \subseteq \mathcal{P}(\mathcal{X})$ be a set of probability measures and $U : \mathbb{R}_{+,\infty} \to \overline{\mathbb{R}}$ a non-decreasing concave $C^1$ function which is bounded from above.
Then the e-variable $f^*$ in Lemma~\ref{lem:exists_integral_utility_eq_sup} is such that for all e-variables $g \in \mathcal{E}_{\mathcal{Q}}$,
\begin{align*}
  P\left[ U'(f^*) (g - f^*) \right] \le 0 \: .
\end{align*}
\end{lemma}

\begin{proof}
For $f \in \mathcal{E}_{\mathcal{Q}}$, let $f_t = t f + (1-t) f^*$ for $t \in [0,1]$.
By concavity of $U$, $U(f_t) \ge t U(f) + (1-t) U(f^*)$, hence
\begin{align*}
  \frac{U(f_t) - U(f^*)}{t}
  &\ge U(f) - U(f^*)
  \: .
\end{align*}
TODO
\end{proof}


\begin{theorem}[Existence of the numeraire]\label{thm:exists_numeraire}
  \uses{def:IsNumeraire}
Let $P \in \mathcal{P}(\mathcal{X})$ and $\mathcal{Q} \subseteq \mathcal{P}(\mathcal{X})$.
A numeraire for $(P, \mathcal{Q})$ exists and is $P$-a.e. unique.
\end{theorem}

\begin{proof}
  \uses{lem:bounded_numeraire_first_order}

\end{proof}

We will write $E^*_{P, \mathcal{Q}}$ for the numeraire for $(P, \mathcal{Q})$.
